\begin{abstract}
Turning a natural-language request into an annealing result requires more than a valid quadratic unconstrained
binary optimization (QUBO) formula. Public data must be acquired and interpreted, the model must preserve the
confirmed requirements, and the delivered decision must be checked independently of its QUBO energy.
We present a human-in-the-loop workflow in which locally served language-model agents clarify requirements,
generate a binary model, acquire and preprocess data, and submit QUBOs to interchangeable classical,
quantum-inspired and hybrid backends. Typed data interfaces, separate data and requirement checkers, and bounded
cross-stage repair provide executable feedback. We also replace reference-bitstring verification with a
problem-level protocol that checks every decoded ground state against a benchmark-anchored evaluator.
An audit found incorrect generated test cases in at least seven of twelve shape failures. Across three runs on
100 benchmarks, problem-level verification yielded 69, 72 and 70 successes; 27 benchmarks changed outcome.
A synthetic 12-category suite produced valid models in all 36 runs. Fixed-model comparisons returned feasible
solutions in all 54 simulated-annealing/GPU conditions and all 27 D-Wave Leap Hybrid conditions, with an explicit
ten-second Hybrid limit. Specified-source real-data pilots establish narrower integration capabilities;
autonomous public-data acquisition and the incremental effects of checkers, repair and retrieval remain to be
evaluated under a frozen repeated protocol. The contribution is a data-grounded workflow and evaluation methodology that distinguishes data and modelling
errors from solver limitations.
\end{abstract}
